% !TEX root = ../Thesis.tex

\chapter{Methodik und Forschungsdesign}
\label{ch:methodik}

% ============================================================================
% HINWEIS: Vorläufiger Inhalt aus dem Exposé
% ============================================================================
\begin{tcolorbox}[
    colback=red!5!white,
    colframe=red!75!black,
    title={\textbf{Hinweis: Vorläufiger Inhalt}},
    fonttitle=\bfseries
]
Die folgenden Abschnitte zur Methodik und zur geplanten Gliederung stammen aus dem Exposé und werden im Laufe der Bachelorarbeit überarbeitet, konkretisiert oder angepasst. Sie dienen derzeit als Arbeitsgrundlage und spiegeln den konzeptionellen Stand zu Beginn der Arbeit wider.
\end{tcolorbox}

\vspace{1em}

Die Arbeit folgt einem artefaktorientierten Forschungsdesign mit konzeptionellem und prototypischem Schwerpunkt. Ziel ist es, eine praxisnahe Observability-Strategie für \emph{frag.jetzt} zu entwickeln und diese anhand konkreter Artefakte nachvollziehbar darzustellen. Die Methodik gliedert sich in sechs aufeinander aufbauende Schritte.

\section{Schritt 1: Service-spezifische Konkretisierung der Four Golden Signals}
\label{sec:methodik-schritt1}

Im ersten Schritt werden die Four Golden Signals als Referenzmodell herangezogen und für die Zielarchitektur von \emph{frag.jetzt} konkretisiert. Für jedes Teilsystem (PWA-Frontend, Spring-Boot-Backend, PostgreSQL-Datenbank und KI-Dienste) werden geeignete Interpretationen für die einzelnen Beobachtungsgrößen abgeleitet, beispielsweise Latenz aus Nutzersicht, Fehlerraten pro Service oder Ressourcenauslastung kritischer Komponenten. Ziel ist die Definition klar interpretierbarer Indikatoren, die sowohl für Dashboards als auch für Alerting nutzbar sind.

\section{Schritt 2: Vergleichende Evaluation von Observability-Stack-Kombinationen}
\label{sec:methodik-schritt2}

Der zweite Schritt umfasst eine vergleichende Evaluation geeigneter Observability-Werkzeuge. Da eine vollständige praktische Erprobung aller verfügbaren Tools im Rahmen der Arbeit nicht realistisch ist, erfolgt der Vergleich primär auf theoretischer und konzeptioneller Ebene. Untersucht werden Prometheus/Grafana, der Elastic Stack sowie Jaeger unter Berücksichtigung der Anforderungen von \emph{frag.jetzt}. Bewertungskriterien sind unter anderem die Abdeckung der benötigten Telemetriearten, die Integrationsfähigkeit in bestehende Systeme, der Betriebsaufwand sowie die Verfügbarkeit als Open-Source-Lösung. Auf dieser Basis wird eine begründete Auswahl eines geeigneten Stacks getroffen.

\section{Schritt 3: Instrumentierung mit OpenTelemetry}
\label{sec:methodik-schritt3}

Im dritten Schritt wird die Instrumentierung der Services konzipiert. OpenTelemetry dient dabei als einheitliches Instrumentierungswerkzeug zur Erfassung von Metriken, Logs und Traces. Ziel ist eine konsistente Telemetrie entlang des gesamten Anfragepfads, die eine Korrelation zwischen Nutzerinteraktionen, Backend-Verarbeitung und externen Abhängigkeiten ermöglicht.

\section{Schritt 4: Rollenbasierte Dashboards}
\label{sec:methodik-schritt4}

Der vierte Schritt befasst sich mit der Konzeption rollenbasierter Dashboards. Für unterschiedliche Zielgruppen wie DevOps, Entwickler und Product Owner werden Sichten entworfen, die Telemetriedaten mit nutzungs- und betriebsrelevantem Kontext verknüpfen. Ziel ist es, nicht nur Auffälligkeiten darzustellen, sondern deren Bedeutung für die jeweilige Rolle nachvollziehbar zu machen, etwa im Hinblick auf Nutzerwahrnehmung, Dringlichkeit oder notwendige Maßnahmen.

\section{Schritt 5: Alerting-Framework mit Fokus auf False-Positive-Reduktion}
\label{sec:methodik-schritt5}

Im fünften Schritt wird ein Alerting-Ansatz entwickelt, der sich an den Golden Signals orientiert und auf Prometheus als technischer Basis aufsetzt. Neben klassischen regelbasierten Alerts werden Konzepte zur Priorisierung und Reduktion irrelevanter Alarme berücksichtigt. Dazu gehören beispielsweise die Aggregation mehrerer Signale, zeitliche Glättung sowie der Einsatz datengetriebener Verfahren wie einfacher Anomalieerkennung. Optional wird untersucht, inwiefern KI-gestützte Verfahren zur Mustererkennung oder Alert-Bewertung eingesetzt werden können, um die Anzahl unnötiger Benachrichtigungen zu reduzieren.

\section{Schritt 6: Entwicklung operativer Runbooks}
\label{sec:methodik-schritt6}

Im sechsten Schritt werden exemplarische Runbooks für wiederkehrende Incident-Szenarien erstellt. Diese Runbooks sollen Alerts mit klaren Diagnose- und Handlungsanweisungen verknüpfen und dienen dazu, Reaktionszeiten zu verkürzen und die operative Qualität der Störungsbearbeitung zu erhöhen.

\section{Ergebnisartefakte}
\label{sec:ergebnisartefakte}

Aus den beschriebenen Schritten ergeben sich die geplanten Ergebnisse:
\begin{enumerate}
    \item Ein wissenschaftlicher Bericht zur Strategie, Toolauswahl und Alerting-Konzeption
    \item Eine lauffähige Observability-Plattform (z.\,B. Prometheus/Grafana)
    \item Rollenbasierte Dashboards
    \item Beispielhafte Runbooks für Incident-Situationen
\end{enumerate}

% ============================================================================
% Geplante Gliederung
% ============================================================================

\section{Geplante Gliederung (konzeptionell)}
\label{sec:geplante-gliederung}

\begin{tcolorbox}[
    colback=red!5!white,
    colframe=red!75!black,
    title={\textbf{Hinweis: Vorläufige Gliederung}},
    fonttitle=\bfseries
]
Die folgende Gliederung ist ein erster Entwurf aus dem Exposé. Die endgültige Struktur der Arbeit kann hiervon abweichen und wird im Verlauf der Bearbeitung angepasst.
\end{tcolorbox}

\vspace{1em}

\begin{enumerate}
    \item \textbf{Einleitung} \\
    Problemstellung, Motivation, Zielsetzung und Forschungsfragen; Einordnung der Arbeit in den Kontext cloud-nativer Systeme und moderner Observability.
    
    \item \textbf{Theoretische Fundierung} \\
    Grundlagen zu Monitoring und Observability (Metriken, Logs, Traces, Alerting); Four Golden Signals als Referenzmodell; Überblick relevanter Observability-Tools und Toolklassen.
    
    \item \textbf{Untersuchungsgegenstand: Architektur von \emph{frag.jetzt}} \\
    Beschreibung der Systemarchitektur (PWA-Frontend, Spring Boot Backend, PostgreSQL, KI-Services sowie Schnittstellen/Kommunikationspfade); Ableitung der Observability-Anforderungen je Komponente.
    
    \item \textbf{Methodik / Forschungsdesign} \\
    Vorgehensmodell der Arbeit; Operationalisierung der Golden Signals; Kriterienkatalog für Tool-/Stack-Evaluation; Instrumentierungs- und Evaluationsplan; Validitäts- und Abgrenzungsaspekte.
    
    \item \textbf{Evaluation von Observability-Stacks} \\
    Vergleich Prometheus/Grafana, ELK Stack, Jaeger (und OpenTelemetry als Instrumentierungsstandard); Ergebnisdarstellung anhand definierter Kriterien; begründete Auswahl einer geeigneten Stack-Kombination.
    
    \item \textbf{Konzeption der Observability-Strategie für \emph{frag.jetzt}} \\
    Service-spezifische Golden-Signal-Definitionen; Messkonzept (Metriken/Logs/Traces) entlang des Anfragepfads; Telemetrie- und Korrelationserfordernisse.
    
    \item \textbf{Umsetzung: Instrumentierung und Observability-Plattform} \\
    Umsetzung der Telemetrie-Erfassung; Konfiguration der Plattformkomponenten; Nachweis der End-to-End-Sicht (Metrik–Trace–Log-Korrelation).
    
    \item \textbf{Dashboards und Alerting-Konzept} \\
    Rollenbasierte Dashboards; Alerting-Regeln auf Grundlage nutzerrelevanter Beobachtungsgrößen; Ansatz zur False-Positive-Reduktion und Priorisierung; Einbindung von Anomalieerkennung (konzeptionell/prototypisch).
    
    \item \textbf{Runbooks und Incident-Playbooks} \\
    Erstellung exemplarischer Runbooks; Verknüpfung von Alerts mit Diagnose- und Handlungsschritten; typische Incident-Szenarien.
    
    \item \textbf{Diskussion} \\
    Einordnung der Ergebnisse, Grenzen der Umsetzung, Trade-offs (z.\,B. Aufwand vs. Nutzen, Datenvolumen/Kardinalität, Abhängigkeit externer KI-Services).
    
    \item \textbf{Fazit und Ausblick} \\
    Zusammenfassung der wichtigsten Ergebnisse; Ausblick auf Weiterentwicklung (z.\,B. Reifegradsteigerung, Service Level Objective (SLO)-basierte Steuerung, weitere Automatisierung).
\end{enumerate}
